% ---------------------------------------------------------------------
%  CHƯƠNG 3 — PHÂN TÍCH VÀ THIẾT KẾ
% ---------------------------------------------------------------------
\chapter{Phân tích và thiết kế}
\label{ch:thietke}

\begin{muctieu}
Chương này phân tích các yêu cầu và mô tả kiến trúc thực tế của phiên bản
VisionAI Pro. Trọng tâm là luồng xử lý ảnh, OCR local, Gemini Vision,
STT/TTS, robot Vision, lệnh thoại, hỏi đáp theo ảnh và lưu lịch sử.
\end{muctieu}

\section{Phân tích bài toán và yêu cầu}

Hệ thống được xây dựng nhằm hỗ trợ người khiếm thị tiếp nhận thông tin từ
hình ảnh và tương tác với trợ lý bằng tiếng Việt. Người dùng có thể tải ảnh,
chụp ảnh chủ động bằng camera, nhập câu hỏi hoặc ghi âm rồi xác nhận gửi.
Ảnh được phân tích bằng Gemini Vision hoặc Tesseract OCR tuỳ chế độ; giọng nói
được chuyển thành văn bản bằng Google Web Speech trước khi gửi câu hỏi tới Gemini.

\subsection{Yêu cầu chức năng}

Các chức năng đã được hiện thực trong source gồm:
\begin{itemize}[leftmargin=1.6em, itemsep=2pt]
  \item Hiển thị giao diện Streamlit với 11 trang: Trang chính, Phân tích hình ảnh, OCR, Hỗ trợ di chuyển, Hỏi đáp AI, Trợ lý giọng nói, Camera, Lịch sử, Thống kê, Cài đặt và Hướng dẫn.
  \item Cho phép tải ảnh JPG/JPEG/PNG/WEBP hoặc chụp ảnh bằng camera.
  \item Phân tích ảnh theo các chế độ: mô tả toàn cảnh, nhận biết vật thể và vị trí, OCR, nhận diện tiền Việt, hỗ trợ người khiếm thị và hỗ trợ di chuyển.
  \item Trả về mô tả, bối cảnh, vật thể, nguy cơ, hướng dẫn di chuyển, văn bản và mệnh giá tiền theo cấu trúc dữ liệu chuẩn hoá.
  \item OCR local bằng Tesseract, không sử dụng credit AI cloud.
  \item Cho phép hỏi tiếp về ảnh gần nhất và duy trì tối đa các lượt hội thoại gần nhất trong phiên.
  \item Nhận giọng nói tiếng Việt bằng SpeechRecognition/Google Web Speech; chỉ xử lý bản ghi sau thao tác xác nhận của người dùng.
  \item Chuyển kết quả văn bản thành giọng nói bằng Edge TTS hoặc Windows SAPI5 qua pyttsx3; hỗ trợ lựa chọn giọng đọc.
  \item Cung cấp robot Vision tương tác, hỏi đáp theo ngữ cảnh trang hiện tại và ảnh gần nhất, tìm vật trong ảnh và lệnh thoại điều hướng/đọc lại.
  \item Lưu lịch sử phân tích vào \texttt{data/history.json}; hỗ trợ lọc, đánh dấu yêu thích, xoá có xác nhận, thống kê và xuất báo cáo TXT.
  \item Cung cấp chế độ trợ năng, bật/tắt giọng đọc, mức chi tiết và lựa chọn giọng trong phiên Streamlit.
\end{itemize}

\subsection{Yêu cầu phi chức năng}

\begin{itemize}[leftmargin=1.6em, itemsep=2pt]
  \item \textbf{Khả năng tiếp cận:} ưu tiên phản hồi âm thanh và nội dung ngắn, dễ nghe.
  \item \textbf{Tính ổn định:} kiểm tra ảnh đầu vào, giới hạn dung lượng, báo lỗi dễ hiểu và thử danh sách model Gemini đã cấu hình.
  \item \textbf{An toàn thông tin:} API key được đọc từ \texttt{.env} và không đưa vào source công khai.
  \item \textbf{Khả năng chạy cục bộ:} giao diện chạy bằng Streamlit trên Windows; OCR dùng Tesseract cục bộ, còn Gemini, Google Web Speech và Edge TTS cần mạng.
  \item \textbf{Khả năng mở rộng:} danh sách model Gemini, ngôn ngữ STT và giọng TTS có thể cấu hình qua biến môi trường.
\end{itemize}

\section{Thiết kế kiến trúc hệ thống}

Kiến trúc thực tế được tổ chức thành bốn nhóm: giao diện Streamlit, xử lý
ảnh/âm thanh, các dịch vụ AI và lớp lưu trữ cục bộ. \texttt{app.py} điều phối giao
diện và gọi các module trong thư mục \texttt{modules/}.

\begin{figure}[H]
\centering
\resizebox{\textwidth}{!}{%
\begin{tikzpicture}[
    font=\small\sffamily,
    box/.style={rectangle, rounded corners=5pt, draw=daunavy, fill=daunavy!7, thick, minimum width=4.2cm, minimum height=0.8cm, align=center},
    ai/.style={rectangle, rounded corners=5pt, draw=daumaroon, fill=daumaroon!7, thick, minimum width=4.2cm, minimum height=0.8cm, align=center},
    data/.style={rectangle, rounded corners=5pt, draw=daunavy, fill=white, thick, minimum width=4.2cm, minimum height=0.8cm, align=center},
    arrow/.style={-Stealth, thick, draw=daunavy}
]
\node (user)     [box]  at (1.25,0)    {Người dùng};
\node (ui)       [box]  at (1.25,-1.4) {Streamlit \texttt{app.py}};
\node (image)    [box]  at (-5,-2.9)   {Ảnh tải lên / camera};
\node (question) [box]  at (5,-2.9)    {Câu hỏi văn bản / audio};

\node (vision)   [ai]   at (-7.5,-4.8) {Gemini Vision\\Google GenAI SDK};
\node (ocr)      [ai]   at (-2.5,-4.8) {Tesseract OCR\\local / chế độ OCR};
\node (stt)      [ai]   at (2.5,-4.8)  {SpeechRecognition\\Google Web Speech};
\node (qa)       [ai]   at (7.5,-4.8)  {Gemini hỏi đáp\\ảnh gần nhất nếu có};

\node (analysis) [box]  at (-5,-6.8)   {Kết quả phân tích ảnh};
\node (answer)   [box]  at (7.5,-6.8)  {Trả lời hỏi đáp};
\node (display)  [box]  at (1.25,-8.6) {Hiển thị trong Streamlit};

\node (history)  [data] at (-9,-8.6)   {history.json\\chỉ lưu kết quả\\phân tích};
\node (tts)      [ai]   at (1.25,-10.4) {Edge TTS\\Windows SAPI5 dự phòng\\khi người dùng yêu cầu đọc};
\node (audio)    [box]  at (1.25,-12.2) {Audio tùy chọn};

\draw [arrow] (user) -- (ui);
\draw [arrow] (ui.south) -- ++(0,-0.4) -| (image.north);
\draw [arrow] (ui.south) -- ++(0,-0.4) -| (question.north);
\draw [arrow] (image.south) -- ++(0,-0.4) -| (vision.north);
\draw [arrow] (image.south) -- ++(0,-0.4) -| (ocr.north);
\draw [arrow] (question.south) -- ++(0,-0.4) -| (stt.north);
\draw [arrow] (question.south) -- ++(0,-0.4) -| (qa.north);
\draw [arrow] (stt.east) -- (qa.west);

\draw [arrow] (vision.south) -- ++(0,-0.5) -| (analysis.north);
\draw [arrow] (ocr.south) -- ++(0,-0.5) -| (analysis.north);
\draw [arrow] (qa.south) -- (answer.north);

\draw [arrow] (analysis.west) -| (history.north);
\draw [arrow] (analysis.south) |- (display.west);
\draw [arrow] (answer.south) |- (display.east);
\draw [arrow] (display) -- (tts);
\draw [arrow] (tts) -- (audio);
\end{tikzpicture}
}
\caption{Kiến trúc thực tế của hệ thống Trợ lý thị giác AI.}
\label{fig:kientruc}
\end{figure}

Hình~\ref{fig:kientruc} tách luồng phân tích ảnh khỏi luồng hỏi đáp. Ảnh được
đưa tới Gemini hoặc Tesseract tuỳ chế độ; khi người dùng đặt câu hỏi bằng
giọng nói, STT chuyển bản ghi thành văn bản trước khi gửi cho Gemini. Câu hỏi
nhập bằng bàn phím bỏ qua STT. TTS chỉ được gọi khi người dùng yêu cầu nghe;
kết quả phân tích ảnh được lưu vào lịch sử cục bộ.

Trong source, OCR local là nhánh riêng với Gemini Vision. STT là một nhánh
riêng nhưng gọi Google Web Speech thông qua
\texttt{SpeechRecognition}, nên cần Internet. Kết quả từ Gemini và OCR đều được
chuẩn hoá để giao diện hiển thị; hỏi đáp có thể kèm ảnh gần nhất làm ngữ cảnh.

\section{Thiết kế luồng xử lý hình ảnh}

Luồng xử lý dưới đây bám sát hàm \texttt{analyze\_image()} và \texttt{analyze\_local\_ocr()}
trong source.

\begin{figure}[H]
\centering
\resizebox{0.95\textwidth}{!}{%
\begin{tikzpicture}[
    font=\small\sffamily,
    block/.style={rectangle, rounded corners=5pt, draw=daunavy, fill=white, thick, minimum width=6.2cm, minimum height=0.8cm, align=center},
    ai/.style={rectangle, rounded corners=5pt, draw=daumaroon, fill=daumaroon!7, thick, minimum width=6.2cm, minimum height=0.8cm, align=center},
    decision/.style={diamond, draw=daunavy, fill=daunavy!7, thick, aspect=2.6, align=center, inner sep=1pt},
    arrow/.style={-Stealth, thick, draw=daunavy},
    lab/.style={font=\scriptsize\sffamily, fill=white, inner sep=1pt}
]
\node (start)    [block]    at (0,0)        {Tải ảnh / Chụp camera};
\node (mode)     [decision] at (0,-1.7)     {Chế độ OCR?};
\node (ocr)      [ai]       at (-4.4,-3.9)  {Chọn OCR:\\Tesseract local\\tiền xử lý + nhận dạng};
\node (vision)   [ai]       at (4.4,-3.9)   {Chế độ khác:\\Gemini Vision\\JSON response};
\node (retry)    [block]    at (4.4,-5.9)   {Thử model theo danh sách\\429: chờ ngắn; 401/403: dừng};
\node (fallback) [block]    at (4.4,-8.0)   {OCR local dự phòng\\nếu tất cả model thất bại\\(chỉ đọc chữ)};
\node (normalize)[block]    at (0,-10.2)    {Chuẩn hoá kết quả};
\node (display)  [block]    at (0,-11.6)    {Hiển thị kết quả + cảnh báo};
\node (save)     [block]    at (-4.4,-13.3) {Lưu history.json\\sau phân tích ảnh};
\node (optional) [block]    at (4.4,-13.3)  {TTS / xuất TXT\\khi người dùng yêu cầu};

\draw [arrow] (start) -- (mode);
\draw [arrow] (mode.west) -| (ocr.north) node[lab, pos=0.25] {Có};
\draw [arrow] (mode.east) -| (vision.north) node[lab, pos=0.25] {Không};
\draw [arrow] (vision) -- (retry);
\draw [arrow] (retry) -- (fallback) node[lab, midway] {thất bại};
\draw [arrow] (ocr.south) |- (normalize.west);
\draw [arrow] (retry.west) -| (normalize.north) node[lab, pos=0.25] {thành công};
\draw [arrow] (fallback.south) |- (normalize.east);
\draw [arrow] (normalize) -- (display);
\draw [arrow] (display.south) -- ++(0,-0.4) -| (save.north);
\draw [arrow] (display.south) -- ++(0,-0.4) -| (optional.north);
\end{tikzpicture}
}
\caption{Sơ đồ luồng xử lý ảnh theo source VisionAI Pro.}
\label{fig:luong_xuly_chitiet}
\end{figure}

Các bước chính gồm:
\begin{enumerate}[leftmargin=1.8em, itemsep=3pt]
  \item Người dùng tải ảnh hoặc chụp ảnh từ camera.
  \item Ứng dụng kiểm tra chế độ. Nếu chọn OCR, ảnh được đưa vào Tesseract local.
  \item Với các chế độ Vision, ảnh được mã hoá thành data URL, tách bytes và MIME type rồi gửi tới Gemini.
  \item Gemini được yêu cầu trả về một JSON object gồm mô tả, độ tin cậy, bối cảnh, vật thể, nguy cơ, navigation, OCR và tiền.
  \item Nếu một model Gemini gặp lỗi, ứng dụng lần lượt thử model tiếp theo; với 429, ứng dụng chờ ngắn trước lần thử tiếp theo, còn 401/403 thì dừng.
  \item Nếu toàn bộ yêu cầu Vision thất bại, source thử OCR local và ghi cảnh báo rõ rằng kết quả OCR không thay thế phân tích cảnh/vật thể.
  \item Kết quả được chuẩn hoá, hiển thị trên giao diện và có thể đọc bằng TTS, lưu lịch sử hoặc xuất TXT.
\end{enumerate}

Giải thuật~\ref{alg:xuly-anh} tóm tắt luồng xử lý trên. Nhánh OCR được chọn
trực tiếp khi người dùng chọn chế độ OCR; OCR cũng được thử làm phương án dự
phòng khi mọi yêu cầu Gemini phát sinh lỗi. Nếu OCR dự phòng không khả dụng,
ứng dụng báo lỗi thay vì tạo kết quả phân tích thay thế.

\begin{algorithm}[H]
\caption{Xử lý ảnh theo chế độ OCR hoặc Vision}
\label{alg:xuly-anh}
\hspace*{\algorithmicindent}\textbf{Đầu vào:} ảnh $I$, chế độ xử lý $m$, danh sách model Gemini $M$\\
\hspace*{\algorithmicindent}\textbf{Đầu ra:} kết quả đã chuẩn hoá và cảnh báo (nếu có)
\begin{algorithmic}[1]
\Function{XuLyAnh}{$I, m, M$}
  \If{$m = \text{OCR}$}
    \State $r \gets \Call{OCRLocal}{I}$
  \Else
    \State $r \gets \Call{GeminiVision}{I, M}$
    \If{Tất cả yêu cầu Gemini phát sinh lỗi}
      \State $r \gets \Call{OCRLocal}{I}$
      \State Gắn cảnh báo rằng OCR không thay thế phân tích cảnh và vật thể
    \EndIf
  \EndIf
  \State $r \gets \Call{ChuanHoa}{r}$
  \State Hiển thị $r$ trên giao diện
  \State Lưu kết quả phân tích ảnh vào history.json
  \State TTS và xuất TXT chỉ thực hiện khi người dùng yêu cầu
  \State \Return $r$
\EndFunction
\end{algorithmic}
\end{algorithm}

\section{Thiết kế luồng hỏi đáp và giọng nói}

\begin{figure}[H]
\centering
\resizebox{0.95\textwidth}{!}{%
\begin{tikzpicture}[
    font=\small\sffamily,
    block/.style={rectangle, rounded corners=5pt, draw=daunavy, fill=white, thick, minimum width=5.2cm, minimum height=0.8cm, align=center},
    ai/.style={rectangle, rounded corners=5pt, draw=daumaroon, fill=daumaroon!7, thick, minimum width=5.2cm, minimum height=0.8cm, align=center},
    decision/.style={diamond, draw=daunavy, fill=daunavy!7, thick, aspect=2.4, align=center, inner sep=1pt},
    arrow/.style={-Stealth, thick, draw=daunavy},
    lab/.style={font=\scriptsize\sffamily, fill=white, inner sep=1pt}
]
\node (text)        [block]    at (-3.6,0)    {Câu hỏi nhập bằng bàn phím};
\node (audio_input) [block]    at (3.6,0)     {Ghi âm + xác nhận\\của người dùng};
\node (stt)         [ai]       at (3.6,-1.9)  {SpeechRecognition\\Google Web Speech};
\node (parser)      [decision] at (0,-4.0)    {Lệnh cục bộ?};
\node (local)       [block]    at (-4.4,-6.4) {Lệnh cục bộ:\\điều hướng / trợ giúp\\đọc lại};
\node (qa)          [ai]       at (4.4,-6.4)  {Hỏi Gemini:\\tìm vật trong ảnh\\ảnh gần nhất nếu có};
\node (answer)      [block]    at (4.4,-8.2)  {Câu trả lời bằng văn bản};
\node (tts)         [ai]       at (4.4,-10.0) {Edge TTS\\Windows SAPI5 dự phòng\\nếu người dùng yêu cầu đọc};
\node (player)      [block]    at (4.4,-11.8) {Audio tùy chọn};

\draw [arrow] (text.south) |- (parser.west);
\draw [arrow] (audio_input) -- (stt);
\draw [arrow] (stt.south) |- (parser.east);
\draw [arrow] (parser.south west) -| (local.north) node[lab, pos=0.25] {Có};
\draw [arrow] (parser.south east) -| (qa.north) node[lab, pos=0.25] {Không};
\draw [arrow] (qa) -- (answer);
\draw [arrow] (answer) -- (tts);
\draw [arrow] (tts) -- (player);
\end{tikzpicture}
}
\caption{Luồng xử lý giọng nói, hỏi đáp theo ảnh và phản hồi âm thanh.}
\label{fig:luong_giongnoi}
\end{figure}

Hình~\ref{fig:luong_giongnoi} mô tả trình tự sau khi người dùng chủ động
ghi âm và xác nhận: nhận dạng trực tuyến, xử lý lệnh thoại hoặc gửi câu hỏi
cùng ngữ cảnh tới Gemini, rồi phát câu trả lời bằng TTS khi được yêu cầu.
Các mũi tên thể hiện thứ tự xử lý, không phải luồng âm thanh thời gian thực.

Người dùng chủ động ghi âm và xác nhận để \texttt{SpeechRecognition} gửi âm
thanh tới Google Web Speech chuyển thành văn bản. Câu hỏi sau đó được gửi cùng
ngữ cảnh trang và ảnh gần nhất (nếu có) tới Gemini. Kết quả có thể được đọc
bằng Edge TTS hoặc Windows SAPI5 khi Edge TTS gặp lỗi. Luồng này không phải
hội thoại âm thanh thời gian thực.

\section{Thiết kế cơ chế cá nhân hoá}

Phiên bản source hiện tại thực hiện cá nhân hoá ở mức phiên làm việc thay vì
xây dựng hồ sơ người dùng hoặc cơ sở dữ liệu. \texttt{st.session\_state} lưu trạng thái
trợ năng, bật/tắt giọng nói, giọng đọc, mức chi tiết, chế độ hiện tại, ảnh gần nhất,
ngữ cảnh trang và lịch sử chat. Ngoài ra, \texttt{history.json} lưu tối đa 100 bản ghi
phân tích để người dùng xem lại, lọc và đánh dấu yêu thích.

\begin{figure}[H]
\centering
\begin{tikzpicture}[
    node distance=0.85cm and 0.9cm,
    font=\small\sffamily,
    box/.style={rectangle, rounded corners=5pt, draw=daunavy, fill=white, thick, minimum width=3.3cm, minimum height=0.72cm, align=center},
    center/.style={rectangle, rounded corners=5pt, draw=daumaroon, fill=daumaroon!7, thick, minimum width=3.6cm, minimum height=0.9cm, align=center},
    arrow/.style={-Stealth, thick, draw=daunavy}
]
\node (settings) [box] {Voice / Detail / Mode};
\node (image) [box, below=of settings] {Ảnh gần nhất};
\node (chat) [box, below=of image] {Ngữ cảnh chat};
\node (session) [center, right=1.3cm of image] {st.session\_state};
\node (analysis) [box, right=1.3cm of session] {Kết quả phân tích ảnh};
\node (history) [box, below=of analysis] {history.json\\tối đa 100 bản ghi};
\node (result) [box, below=1.0cm of session] {Kết quả và trải nghiệm phiên};
\draw [arrow] (settings) -- (session);
\draw [arrow] (image) -- (session);
\draw [arrow] (chat) -- (session);
\draw [arrow] (session) -- (result);
\draw [arrow] (session) -- (analysis);
\draw [arrow] (analysis) -- (history);
\end{tikzpicture}
\caption{Cơ chế cá nhân hoá và lưu trạng thái thực tế của phiên bản hiện tại.}
\label{fig:canhanhoa}
\end{figure}

Hình~\ref{fig:canhanhoa} thể hiện các tuỳ chọn và ngữ cảnh tạm thời được giữ
trong phiên Streamlit, còn các kết quả được lưu riêng trong tệp lịch sử.
Điều này giải thích vì sao phiên bản hiện tại chưa có hồ sơ cá nhân lâu dài.

Cách tiếp cận này phù hợp với phạm vi nguyên mẫu. Source chưa có đăng nhập,
hồ sơ người dùng hoặc cơ sở dữ liệu quan hệ; vì vậy báo cáo không mô tả các
thành phần này như một chức năng đã triển khai.

\section{Thiết kế dữ liệu}

Dữ liệu phân tích được biểu diễn bằng cấu trúc JSON gồm các trường \texttt{summary},
\texttt{confidence}, \texttt{scene}, \texttt{objects}, \texttt{accessibility\_tips}, \texttt{text\_detected},
\texttt{money\_detected}, \texttt{hazards}, \texttt{navigation} và \texttt{warning}. Lịch sử được lưu trong
\texttt{data/history.json} và mỗi lần lưu chỉ giữ tối đa 100 bản ghi gần nhất.

\section{Thiết kế giao diện}

Giao diện được xây dựng bằng Streamlit với thanh điều hướng bên trái và một
robot Vision đồ hoạ có thể mở chat hoặc điều hướng tới các công cụ. Các trang
gồm trang chính, phân tích ảnh, OCR, hỗ trợ di chuyển, hỏi đáp AI, trợ lý giọng
nói, camera, lịch sử, thống kê, cài đặt và hướng dẫn. Chế độ trợ năng tăng cỡ
chữ/nút và làm rõ focus bàn phím; người dùng có thể đổi giọng đọc, bật/tắt
giọng nói, kiểm tra kết nối Gemini và xoá phiên hoặc lịch sử.

Hình~\ref{fig:bocuc-web} mô tả bố cục trang chính. Thanh bên chứa điều hướng
11 trang cùng các công tắc trợ năng; khu vực chính gồm robot Vision, hàng
thẻ công cụ, kết quả gần nhất và cột hoạt động gần đây.

\begin{figure}[H]
\centering
\resizebox{\textwidth}{!}{%
\begin{tikzpicture}[
    font=\scriptsize\sffamily,
    frame/.style={rectangle, rounded corners=4pt, draw=daunavy, thick, fill=white},
    side/.style={rectangle, rounded corners=4pt, draw=daunavy, thick, fill=daunavy!7},
    hero/.style={rectangle, rounded corners=4pt, draw=daumaroon, thick, fill=daumaroon!7},
    card/.style={rectangle, rounded corners=3pt, draw=daunavy, thick, fill=daunavy!7, align=center, minimum width=1.38cm, minimum height=1.3cm, inner sep=2pt},
    btn/.style={rectangle, rounded corners=2pt, draw=daunavy, fill=white, align=center, inner sep=2pt, minimum height=0.42cm},
    lbl/.style={font=\scriptsize\sffamily\bfseries, text=daunavy, anchor=west}
]
% khung ngoài
\draw[frame] (0,0) rectangle (15.4,8.9);
% ---------- thanh bên ----------
\draw[side] (0.15,0.15) rectangle (3.75,8.75);
\node[anchor=north west, font=\scriptsize\sffamily\bfseries, text=daumaroon] at (0.3,8.6) {VisionAI Pro};
\node[anchor=north west, font=\tiny\sffamily] at (0.3,8.25) {Trạng thái AI: Đã sẵn sàng};
\node[anchor=north west, align=left, font=\tiny\sffamily, text width=3.2cm] at (0.3,7.85)
  {Điều hướng (11 trang):\\[1pt]
   Trang chính\\ Phân tích hình ảnh\\ OCR\\ Hỗ trợ di chuyển\\ Hỏi đáp AI\\
   Trợ lý giọng nói\\ Camera\\ Lịch sử\\ Thống kê\\ Cài đặt\\ Hướng dẫn};
\node[btn, minimum width=3.2cm] at (1.95,3.05) {Trò chuyện với Vision};
\node[btn, minimum width=3.2cm] at (1.95,2.4) {Công tắc giọng nói};
\node[btn, minimum width=3.2cm] at (1.95,1.75) {Công tắc trợ năng};
\node[btn, minimum width=3.2cm] at (1.95,0.9) {Xóa phiên hiện tại};
% ---------- robot ----------
\draw[hero] (3.95,6.55) rectangle (15.25,8.75);
\node[btn, minimum width=4.2cm, align=left] at (6.15,8.2) {Lời thoại của robot Vision};
\node[anchor=west, font=\small\sffamily\bfseries] at (4.15,7.45) {VisionAI Pro};
\node[btn, anchor=west] (hb1) at (4.15,6.9) {Phân tích hình ảnh};
\node[btn, anchor=west] (hb2) at ([xshift=0.2cm]hb1.east) {Đọc văn bản};
\node[btn, anchor=west] at ([xshift=0.2cm]hb2.east) {Hỏi đáp theo ảnh};
\node[rectangle, rounded corners=6pt, draw=daumaroon, fill=white, minimum width=2.4cm, minimum height=1.7cm, align=center] at (13.3,7.65) {Robot Vision\\(16 trạng thái\\biểu cảm)};
% ---------- cột chính ----------
\node[lbl] at (3.95,6.2) {Công cụ};
\foreach \i/\t in {0/{Phân tích\\ảnh},1/{Đọc\\văn bản},2/{Hỏi đáp\\ảnh},3/{Giọng\\nói},4/{Camera}} {
  \node[card] at (4.65+\i*1.52,5.35) {\t};
  \node[btn, minimum width=1.38cm, font=\tiny\sffamily] at (4.65+\i*1.52,4.45) {Mở};
}
\node[lbl] at (3.95,3.85) {Kết quả gần nhất};
\draw[frame] (3.95,0.3) rectangle (11.55,3.6);
\node[align=left, anchor=north west, font=\tiny\sffamily] at (4.1,3.45)
  {Mô tả, bối cảnh, cảnh báo an toàn, vật thể,\\
   tiền Việt, văn bản OCR, gợi ý cho người khiếm thị,\\
   nút nghe kết quả và nút xuất báo cáo TXT.};
% ---------- cột phải ----------
\node[btn, minimum width=3.4cm] at (13.4,6.2) {Cấu hình Gemini};
\draw[frame] (11.75,0.3) rectangle (15.25,5.85);
\node[anchor=north west, font=\scriptsize\sffamily\bfseries, text=daunavy] at (11.9,5.7) {Hoạt động gần đây};
\node[anchor=north west, align=left, font=\tiny\sffamily] at (11.9,5.25)
  {Tối đa 4 bản ghi\\từ \texttt{history.json}\\(chế độ và thời điểm\\phân tích).};
\end{tikzpicture}}
\caption[Sơ đồ bố cục trang chính của giao diện web]{Sơ đồ bố cục trang chính của giao diện web, vẽ lại từ cấu trúc trong \texttt{app.py}; đây là sơ đồ thiết kế, không phải ảnh chụp màn hình.}
\label{fig:bocuc-web}
\end{figure}

\section{Thiết kế thành phần AI}

Thành phần AI được tách thành \texttt{modules/ai\_vision.py}, \texttt{modules/local\_ocr.py}
và \texttt{modules/speech.py}. \texttt{ai\_vision.py} phụ trách Gemini Vision và hỏi đáp theo
ảnh; \texttt{local\_ocr.py} phụ trách Tesseract; \texttt{speech.py} phụ trách STT/TTS. Cách
tách module giúp thay đổi từng thành phần mà không phải sửa toàn bộ giao diện.